\documentclass[11pt,a4paper]{article}
\usepackage[utf8]{inputenc}
\usepackage[spanish]{babel}
\usepackage{amsmath}
\usepackage{amsfonts}
\usepackage{amssymb}
\usepackage{booktabs}
\usepackage{geometry}
\usepackage{hyperref}
\usepackage{titlesec}
\usepackage{xcolor}

% Configuración de márgenes
\geometry{top=2.5cm, bottom=2.5cm, left=3.0cm, right=3.0cm}

% Configuración de hipervínculos
\hypersetup{
    colorlinks=true,
    linkcolor=blue,
    filecolor=magenta,      
    urlcolor=cyan,
    pdftitle={Informe de Analisis Econometrico de Portafolios},
    pdfpagemode=FullScreen,
}

% Formato de secciones
\titleformat{\section}
  {\normalfont\Large\bfseries\color{blue!80!black}}
  {\thesection}{1em}{}
\titleformat{\subsection}
  {\normalfont\large\bfseries\color{blue!60!black}}
  {\thesubsection}{1em}{}

\title{\textbf{Informe de Análisis Econométrico y Co-movimiento de Portafolios}}
\author{\textbf{Lead-Flow Inmobiliaria} \\ Analista Cuantitativo Senior}
\date{1 de julio de 2026}

\begin{document}

\maketitle

\begin{abstract}
Este informe presenta el análisis econométrico cuantitativo de la evolución de dos portafolios de inversión administrados durante el período comprendido entre el 15/02/2022 y el 16/02/2023. Investigamos la aparente paradoja estadística de una correlación diaria extremadamente alta ($r = 0.993$) frente a una ausencia completa de cointegración ($p = 0.730$). Aportamos justificaciones teóricas basadas en la presencia de raíces unitarias (tendencias estocásticas), la regresión espuria y las diferencias estructurales de asignación de activos en renta fija en el contexto macroeconómico del ciclo alcista de tasas de la Reserva Federal durante el año 2022.
\end{abstract}

\newpage

\section{Introducción y Planteamiento}
En la práctica del análisis de portafolios y gestión de riesgos, es común evaluar el grado de co-movimiento entre distintas carteras para estimar la diversificación y la sensibilidad sistemática. Sin embargo, los coeficientes de correlación clásicos pueden resultar engañosos al tratar con series de tiempo financieras de valorización absoluta. 

Este estudio examina la relación temporal entre el \textbf{Portafolio 1} ($V_{1,t}$) y el \textbf{Portafolio 2} ($V_{2,t}$). A través de contrastes de raíz unitaria (ADF) y de cointegración de Engle-Granger, demostramos que la aparente sincronía de largo plazo es de carácter transitorio, y que ambas carteras poseen trayectorias económetricamente independientes y divergentes en el largo plazo.

\section{Análisis de Co-movimiento y Estructura}
Al analizar las $367$ observaciones diarias de valorización de ambos portafolios, se obtuvieron los parámetros estadísticos descritos en la Tabla~\ref{tab:correlacion}.

\begin{table}[h]
\centering
\caption{Parámetros de Co-movimiento de Corto Plazo}
\label{tab:correlacion}
\vspace{0.2cm}
\begin{tabular}{lc}
\toprule
\textbf{Estadístico Cuantitativo} & \textbf{Valor Observado} \\
\midrule
Correlación en Niveles ($V_{1,t}$ vs $V_{2,t}$) & $0.995447$ \\
Correlación en Retornos Diarios ($R_{1,t}$ vs $R_{2,t}$) & $0.992968$ \\
Sensibilidad ($\beta$ de $P_1$ frente a $P_2$) & $1.031237$ \\
Coeficiente de Determinación ($R^2$) & $98.60\%$ \\
\bottomrule
\end{tabular}
\end{table}

\subsection{Solapamiento de Activos en Común}
La altísima correlación en los retornos diarios del $99.30\%$ se debe a un solapamiento estructural crítico en la asignación de activos al inicio del período ($t=0$):
\begin{enumerate}
    \item \textbf{Asignación Idéntica:} De los 17 activos disponibles, 8 de ellos comparten exactamente el mismo peso, sumando un \textbf{$43.3\%$} de la asignación total de las carteras. El activo más grande de ambos portafolios es la renta variable de EEUU, con un peso idéntico del \textbf{$28.0\%$}.
    \item \textbf{Estructura Balanceada:} Ambos portafolios se estructuran bajo un perfil balanceado clásico. El Portafolio 1 mantiene un $52\%$ en renta variable frente a un $50\%$ del Portafolio 2.
\end{enumerate}

La única diferencia estructural radica en la composición de la porción de renta fija ($17\%$ del portafolio): el Portafolio 1 busca mayor rendimiento neto mediante la asignación en \textit{High Yield} y \textit{Emerging Markets Debt}, mientras que el Portafolio 2 desplaza ese $17\%$ hacia activos de alta calidad crediticia y refugio (\textit{US Treasuries}, \textit{MBS} e \textit{Investment Grade Corporates}).

\section{Análisis de Raíz Unitaria (Test ADF)}
Para verificar si las series de valorización del portafolio poseen una tendencia de carácter determinista o estocástica, ejecutamos el test de Dickey-Fuller Aumentado (ADF) con constante y tendencia temporal sobre la serie en niveles de cada portafolio.

\subsection{Especificación del Modelo}
El modelo de regresión estimado para el test ADF está definido por:
\begin{equation}
    \Delta V_t = \alpha + \beta t + \gamma V_{t-1} + \sum_{j=1}^{p} \delta_j \Delta V_{t-j} + \epsilon_t
\end{equation}
Donde las hipótesis evaluadas son:
\begin{itemize}
    \item $H_0: \gamma = 0$ (Presencia de Raíz Unitaria $\rightarrow$ Tendencia Estocástica).
    \item $H_1: \gamma < 0$ (Estacionario en Tendencia $\rightarrow$ Tendencia Determinista).
\end{itemize}

\subsection{Resultados Empíricos}
Los resultados obtenidos para el Portafolio 1 se detallan en la Tabla~\ref{tab:adf}.

\begin{table}[h]
\centering
\caption{Resultados del Test ADF (Portafolio 1)}
\label{tab:adf}
\vspace{0.2cm}
\begin{tabular}{lc}
\toprule
\textbf{Parámetro del Contraste} & \textbf{Valor / Estadístico} \\
\midrule
Estadístico del Test ($\tau_{\text{ADF}}$) & $-1.63204$ \\
p-valor (Valor de Probabilidad) & $0.77964$ \\
Valor Crítico (1\%) & $-3.98393$ \\
Valor Crítico (5\%) & $-3.42265$ \\
Valor Crítico (10\%) & $-3.13420$ \\
\bottomrule
\end{tabular}
\end{table}

\textbf{Conclusión Estadística:} Dado que el estadístico calculado ($\tau_{\text{ADF}} = -1.63204$) es superior (menos negativo) al valor crítico al $5\%$ ($-3.42265$), \textbf{no se puede rechazar la hipótesis nula $H_0$}. La valoración del portafolio posee una tendencia estocástica (raíz unitaria). Los shocks y fluctuaciones del mercado se acumulan de forma permanente en la trayectoria del portafolio.

\section{Análisis de Cointegración (Test de Engle-Granger)}
Para comprobar si existe una relación de equilibrio o tendencia común de largo plazo que vincule las valorizaciones de ambos portafolios, ejecutamos el contraste de cointegración de Engle-Granger de dos pasos.

\subsection{Especificación del Modelo}
En el primer paso, corremos la regresión de cointegración mediante MCO:
\begin{equation}
    V_{1,t} = \mu + \theta V_{2,t} + u_t
\end{equation}
En el segundo paso, se contrasta la presencia de raíz unitaria sobre la serie de residuos estimados $\hat{u}_t$ mediante una prueba ADF:
\begin{equation}
    \Delta \hat{u}_t = \gamma \hat{u}_{t-1} + \sum_{j=1}^{p} \delta_j \Delta \hat{u}_{t-j} + e_t
\end{equation}
Donde la hipótesis nula $H_0: \gamma = 0$ postula la ausencia de cointegración (el spread se desvía permanentemente y diverge).

\subsection{Resultados Empíricos}
Los estadísticos de la prueba de residuos se resumen en la Tabla~\ref{tab:coint}.

\begin{table}[h]
\centering
\caption{Resultados de la Cointegración de Engle-Granger}
\label{tab:coint}
\vspace{0.2cm}
\begin{tabular}{lc}
\toprule
\textbf{Estadístico de Contraste} & \textbf{Valor Observado} \\
\midrule
Estadístico de Cointegración ($\tau_{\text{coint}}$) & $-2.09385$ \\
p-valor de Cointegración & $0.73075$ \\
Relación de Largo Plazo & \textbf{Divergente (Sin Cointegración)} \\
\bottomrule
\end{tabular}
\end{table}

\textbf{Conclusión Estadística:} El estadístico de cointegración ($\tau_{\text{coint}} = -2.09385$) es incapaz de superar los límites críticos de Engle-Granger. Por tanto, \textbf{no rechazamos la hipótesis nula de ausencia de cointegración}. El spread entre ambos portafolios es una serie integrada de orden 1 ($I(1)$), lo que significa que los portafolios pueden divergir de forma permanente.

\section{Interpretación Económica y Financiera}
La aparente contradicción entre la correlación del $99.30\%$ y el rechazo de la cointegración ilustra perfectamente la diferencia entre los comovimientos de corto plazo y el equilibrio estructural de largo plazo:

\begin{enumerate}
    \item \textbf{El Shock Sistémico de 2022:}
    Durante el período analizado, el mercado mundial experimentó un choque común de carácter macroeconómico: la subida masiva de tasas por parte de la Reserva Federal. Esto provocó una corrección simultánea en activos de renta variable y fija globales, haciendo caer a ambas carteras al unísono (comovimiento a corto plazo).
    \item \textbf{La Regresión Espuria:}
    La presencia de raíces unitarias en ambas series ($I(1)$), combinada con una tendencia macro común, genera un sesgo de regresión espuria en los estadísticos clásicos. La correlación en niveles ($0.995$) es puramente estadística y no representa un lazo de equilibrio real.
    \item \textbf{Factores de Riesgo Diferenciados:}
    El $17\%$ de asignación diferencial (High Yield y Deuda Emergente en $P_1$ vs. Bonos Soberanos en $P_2$) tiene dinámicas e incentivos económicos propios. Ante un escenario de estrés crediticio o recesión prolongada, la probabilidad de impago de la cartera de High Yield generará pérdidas asimétricas que desligarán permanentemente al Portafolio 1 del Portafolio 2, confirmando que las carteras son fundamentalmente divergentes.
\end{enumerate}

\section{Conclusiones}
El análisis econométrico demuestra que la correlación histórica es un indicador insuficiente para modelar el equilibrio de largo plazo entre portafolios. El contraste formal de Engle-Granger confirma que las diferencias en la calidad crediticia de la porción de renta fija de ambas carteras introduce una deriva estocástica independiente. Las carteras no comparten una tendencia común en el sentido cuantitativo de cointegración, por lo que su gestión y control de riesgo deben abordarse de manera independiente.

\end{document}
